%% ----------------------------------------------------------------------------
%%  §5.5  Derivate delle funzioni elementari
%% ----------------------------------------------------------------------------
\section{Derivate delle funzioni elementari}
\label{sec:derivate-elementari}

In questo paragrafo calcoliamo le derivate delle funzioni elementari. Insieme alle regole
dei due paragrafi precedenti, esse permettono di derivare tutte le funzioni che si
ottengono da quelle elementari mediante somme, prodotti, quozienti, composizioni e
inversioni.

\subsection{Potenze con esponente naturale e polinomi}

\begin{proposizione}\label{prop:derivata-xn}
Per ogni $n\in\N$ e ogni $x\in\R$ si ha
\begin{equation}\label{eq:derivata-xn}
  Dx^{n}=nx^{n-1}
\end{equation}
(per $n=1$ si intende $Dx=1$).
\end{proposizione}

\begin{schema}
\Ipotesi{$n\in\N$; \ $x\in\R$.}
\Tesi{$x^{n}$ è derivabile in $x$ e $Dx^{n}=nx^{n-1}$.}
\end{schema}

\begin{strategia}
La formula dipende da un indice $n\in\N$: è naturale procedere per induzione. Il caso
$n=1$ è la derivata della funzione lineare $x$, già calcolata nell'\cref{es:derivata-lineare};
nel passo induttivo si scrive $x^{n+1}$ come prodotto $x^{n}\cdot x$ e si applica la regola
del prodotto \eqref{eq:derivata-prodotto}, usando l'ipotesi induttiva per la derivata del
primo fattore.
\end{strategia}

\begin{dimostrazione}
Questa formula (che per $x>0$ si ottiene anche dalla più generale \eqref{eq:derivata-xb}) si può
dimostrare facendo uso del principio di induzione (\cref{teo:induzione}): abbiamo già
verificato che $Dx=1$ \eqref{eq:derivata-retta}; quindi la \eqref{eq:derivata-xn} è vera
per $n=1$. Supponiamo, secondo lo schema del principio di induzione, che la
\eqref{eq:derivata-xn} sia vera per un certo $n$, e calcoliamo per mezzo della regola di
derivazione del prodotto:
\begin{equation}\label{eq:derivata-xn-passo}
\begin{split}
  Dx^{n+1}&=D(x^{n}\cdot x)=D(x^{n})\,x+x^{n}\,Dx\\
  &=nx^{n-1}\,x+x^{n}\cdot1=(n+1)x^{n}.
\end{split}
\end{equation}
Abbiamo quindi verificato che la \eqref{eq:derivata-xn} vale anche per l'indice $n+1$.
Perciò la \eqref{eq:derivata-xn} è provata per ogni $n\in\N$.
\end{dimostrazione}

\begin{strategia}[Un'altra via]
Al posto dell'induzione, una formula chiusa, come in Giusti \cite{giusti}. La
scomposizione \eqref{eq:scomposizione-potenze} di $u^{n}-v^{n}$ permette di semplificare
il fattore $\xi-x$ nel rapporto incrementale \eqref{eq:derivata-xi}, che diventa una
somma di $n$ monomi; ciascuno tende a $x^{n-1}$, e il coefficiente $n$ della formula
conta proprio questi addendi. La regola del prodotto non serve.
\end{strategia}

\begin{dimostrazione}[Seconda dimostrazione]
La \eqref{eq:scomposizione-potenze}, ricavata nel \cref{lem:stima-potenze} per
$0\le v<u$, vale per ogni coppia di numeri reali $u$, $v$ (con la convenzione $t^{0}=1$
per ogni $t$). Detta infatti $S=u^{n-1}+u^{n-2}v+\dots+uv^{n-2}+v^{n-1}$ la somma tra
parentesi, si ha
\[
  uS=u^{n}+u^{n-1}v+\dots+uv^{n-1}, \qquad vS=u^{n-1}v+\dots+uv^{n-1}+v^{n},
\]
e nella differenza $uS-vS$ tutti i termini si cancellano, tranne $u^{n}-v^{n}$. Con
$u=\xi$, $v=x$ e $\xi\neq x$ si ottiene
\[
  \frac{\xi^{n}-x^{n}}{\xi-x}=\xi^{n-1}+\xi^{n-2}x+\dots+\xi x^{n-2}+x^{n-1}
  =\sum_{k=0}^{n-1}\xi^{n-1-k}x^{k}.
\]
Per $\xi\to x$ ciascuno degli $n$ addendi tende a $x^{n-1-k}x^{k}=x^{n-1}$, per il limite
di un prodotto (\cref{teo:operazioni-funzioni}); quindi il rapporto incrementale tende a
$nx^{n-1}$, e per la \eqref{eq:derivata-xi} si ha $Dx^{n}=nx^{n-1}$.
\end{dimostrazione}

Notiamo che il risultato ottenuto, insieme alla linearità della derivata, ci permette di
calcolare la derivata di un \emph{polinomio} qualsiasi:
\begin{equation}\label{eq:derivata-polinomio}
\begin{gathered}
  y=a_{n}x^{n}+a_{n-1}x^{n-1}+\dots+a_{1}x+a_{0};\\
  y'=na_{n}x^{n-1}+(n-1)a_{n-1}x^{n-2}+\dots+a_{1}.
\end{gathered}
\end{equation}
In particolare, la derivata di un polinomio di grado $n\in\N$ è un polinomio di grado
$n-1$, e un polinomio è derivabile infinite volte (è di classe $C^{\infty}$). Per gli
esponenti interi negativi la regola del reciproco \eqref{eq:derivata-reciproco} dà, per
$x\neq0$, $Dx^{-n}=D\frac1{x^{n}}=-\frac{nx^{n-1}}{x^{2n}}=-nx^{-n-1}$: la
\eqref{eq:derivata-xn} vale per ogni esponente intero.

\subsection{Logaritmi ed esponenziali}

\begin{proposizione}[Derivata del logaritmo]\label{prop:derivata-log}
Per ogni base $a>0$, $a\neq1$, e per ogni $x>0$ si ha
\begin{equation}\label{eq:derivata-loga}
  D\log_{a}x=\frac1x\log_{a}e=\frac1{x\log a};
\end{equation}
in particolare, per il logaritmo naturale,
\begin{equation}\label{eq:derivata-log}
  D\log x=\frac1x \qquad\perogni x>0 .
\end{equation}
\end{proposizione}

\begin{schema}
\Ipotesi{$a>0$, \ $a\neq1$; \ $x>0$.}
\Tesi{$\log_{a}$ è derivabile in $x$ e $D\log_{a}x=\frac1{x\log a}$; \ in particolare
$D\log x=\frac1x$.}
\end{schema}

\begin{strategia}
Le proprietà del logaritmo trasformano la differenza $\log(x+h)-\log x$ nel logaritmo di
un rapporto, $\log\bigl(1+\frac hx\bigr)$; il rapporto incrementale diventa allora,
a meno del fattore $\frac1x$, la quantità $\frac{\log(1+t)}{t}$ con $t=\frac hx\to0$, di cui
conosciamo il limite dal \cref{cap:funzioni-continue}. È qui che si vede perché i
matematici preferiscono la base $e$: è l'unica per cui quel limite vale $1$.
\end{strategia}

\begin{dimostrazione}
Sia $x>0$ e sia $0<\abs h<x$, così che $x+h>0$. Per la \eqref{eq:log2}
\begin{equation}\label{eq:rapporto-log}
  \frac{\log(x+h)-\log x}{h}=\frac1h\log\frac{x+h}{x}
  =\frac1x\cdot\frac{\log\bigl(1+\frac hx\bigr)}{\frac hx}.
\end{equation}
La funzione $t=g(h)=\frac hx$ tende a $0$ per $h\to0$ ed è diversa da zero per $h\neq0$;
per il limite notevole \eqref{eq:ln-log} e per il \cref{teo:limite-composta}, l'ultimo
fattore tende a $1$. Dunque il rapporto incrementale tende a $\frac1x$, cioè vale la
\eqref{eq:derivata-log}. Per una base qualsiasi si ha $\log_{a}x=\frac{\log x}{\log a}$
\eqref{eq:log4}, formula valida anche per esponenti reali (\cref{sec:monotone-inverse}),
e quindi $D\log_{a}x=\frac1{x\log a}$; infine $\frac1{\log a}=\log_{a}e$, sempre per la
\eqref{eq:log4}.
\end{dimostrazione}

Risulta ora chiaro l'interesse nel considerare logaritmi in base $e$: dato che
$\log_{e}e=1$, la derivata del logaritmo in base $e$ di $x$ è semplicemente $\frac1x$.

\begin{proposizione}[Derivata dell'esponenziale]\label{prop:derivata-exp}
Per ogni $x\in\R$ e ogni $a>0$ si ha
\begin{equation}\label{eq:derivata-exp}
  De^{x}=e^{x}, \qquad Da^{x}=a^{x}\log a .
\end{equation}
\end{proposizione}

\begin{schema}
\Ipotesi{$x\in\R$; \ $a>0$.}
\Tesi{$e^{x}$ e $a^{x}$ sono derivabili in $x$, \ $De^{x}=e^{x}$, \ $Da^{x}=a^{x}\log a$.}
\end{schema}

\begin{strategia}
L'esponenziale è la funzione inversa del logaritmo, di cui ora conosciamo la derivata
$\frac1x$, che non si annulla mai: il teorema di derivazione delle funzioni inverse
(\cref{teo:derivata-inversa}) dà la derivata di $e^{y}$ come reciproco di $\frac1x$,
calcolato nel punto $x=e^{y}$. Per una base qualsiasi si scrive $a^{x}=e^{x\log a}$ e si
applica la regola della catena. Dopo la dimostrazione indicheremo anche una via diretta,
basata sul limite notevole di $\frac{e^{h}-1}{h}$.
\end{strategia}

\begin{dimostrazione}
La funzione $y=\log x$ è continua, strettamente crescente e invertibile, e la sua
inversa è $x=e^{y}$. Fissato $y$, sia $x=e^{y}$: poiché $D\log x=\frac1x\neq0$, il
teorema di derivazione delle funzioni inverse (applicato al logaritmo su un intervallo
chiuso contenente $x$ al suo interno) dà
\begin{equation}\label{eq:derivata-exp-inversa}
  De^{y}=\frac1{D\log x}=\frac1{1/x}=x=e^{y}.
\end{equation}
Usando, come si è soliti fare, il simbolo $x$ per denotare la variabile indipendente,
possiamo riscrivere la formula precedente: $De^{x}=e^{x}$. Per una base $a>0$ qualsiasi
ricordiamo che $a^{x}=e^{x\log a}$ \eqref{eq:potenza-esponenziale}; facendo uso del
teorema di derivazione delle funzioni composte otteniamo
\begin{equation}\label{eq:derivata-ax}
  Da^{x}=De^{x\log a}=e^{x\log a}\,D(x\log a)=a^{x}\log a . \qedhere
\end{equation}
\end{dimostrazione}

La formula $De^{x}=e^{x}$ si può ottenere anche direttamente dalla definizione:
$\frac{e^{x+h}-e^{x}}{h}=e^{x}\cdot\frac{e^{h}-1}{h}\to e^{x}\cdot1$, per il limite
notevole \eqref{eq:ln-exp}. Essa esprime una proprietà caratteristica dell'esponenziale
naturale: \emph{il tasso di accrescimento di $e^{x}$ è uguale al suo valore}. È per
questo che la funzione $e^{x}$ descrive tutti i fenomeni in cui la rapidità di
variazione di una grandezza è proporzionale alla grandezza stessa: la crescita di una
popolazione, il decadimento radioattivo, la scarica di un condensatore, e, come vedremo
studiando le equazioni differenziali, la risposta di molti sistemi dinamici.

\subsection{Potenze con esponente reale}

Ricordiamo la proprietà $e^{\log x}=x$ \eqref{eq:log-inversa}, che è sempre utile quando
si vuole calcolare la derivata di un esponenziale o di una potenza che non rientrano nei
casi precedenti. Analogamente a quanto fatto per $a^{x}$, si calcola la derivata della
funzione potenza $x^{b}$, con esponente $b$ \emph{reale}: per $x>0$, essendo
$x^{b}=e^{b\log x}$ \eqref{eq:potenza-esponenziale},
\begin{equation}\label{eq:derivata-xb}
  Dx^{b}=De^{b\log x}=e^{b\log x}\,D(b\log x)=x^{b}\cdot\frac bx=bx^{b-1}.
\end{equation}
La formula precedente è molto utile. È utilizzata ad esempio nei casi $b=\frac12$ (in
questo caso si riottiene la \eqref{eq:derivata-radice}) e $b=-1$:
\begin{equation}\label{eq:derivata-radice-reciproco}
\begin{gathered}
  D\sqrt x=Dx^{1/2}=\frac12x^{-1/2}=\frac1{2\sqrt x},\\
  D\frac1x=Dx^{-1}=(-1)x^{-2}=-\frac1{x^{2}} .
\end{gathered}
\end{equation}
Lo stesso artificio permette di derivare le funzioni della forma $f(x)^{g(x)}$ con
$f(x)>0$, scrivendole come $e^{g(x)\log f(x)}$. Ad esempio, per $x>0$,
\begin{equation}\label{eq:derivata-x-alla-x}
  Dx^{x}=De^{x\log x}=e^{x\log x}\Bigl(\log x+x\cdot\frac1x\Bigr)=x^{x}(\log x+1).
\end{equation}

\subsection{Funzioni trigonometriche}

\begin{proposizione}[Derivate di seno, coseno e tangente]\label{prop:derivata-sen-cos}
Per ogni $x\in\R$ si ha
\begin{equation}\label{eq:derivata-sen-cos}
  D\sen x=\cos x, \qquad D\cos x=-\sen x,
\end{equation}
e per ogni $x\neq\frac\pi2+k\pi$, $k\in\Z$,
\begin{equation}\label{eq:derivata-tg}
  D\tg x=\frac1{\cos^{2}x}=1+\tg^{2}x .
\end{equation}
\end{proposizione}

\begin{schema}
\Ipotesi{$x\in\R$ (angoli in radianti); \ per la tangente, $\cos x\neq0$.}
\Tesi{$D\sen x=\cos x$; \ $D\cos x=-\sen x$; \ $D\tg x=1+\tg^{2}x$.}
\end{schema}

\begin{strategia}
Per il seno e il coseno si scrive il rapporto incrementale, si sviluppano $\sen(x+h)$ e
$\cos(x+h)$ con le formule di addizione e si raccolgono i termini in modo che $h$ compaia
soltanto nei due limiti notevoli $\frac{\sen h}{h}\to1$ e $\frac{1-\cos h}{h}\to0$; i
fattori $\sen x$ e $\cos x$, che non dipendono da $h$, escono dal limite. La tangente è il
quoziente $\frac{\sen x}{\cos x}$: basta la regola del quoziente, insieme alla relazione
fondamentale $\sen^{2}x+\cos^{2}x=1$.
\end{strategia}

\begin{dimostrazione}
Dimostriamo la prima delle \eqref{eq:derivata-sen-cos}: facciamo uso della formula di
addizione \eqref{eq:addizione-sen} e dei limiti notevoli \eqref{eq:lim-senx-x},
\eqref{eq:lim-1-cos-x}:
\begin{equation}\label{eq:derivata-sen}
\begin{split}
  \lim_{h\to0}\frac{\sen(x+h)-\sen x}{h}
  &=\lim_{h\to0}\frac{\sen x\cos h+\sen h\cos x-\sen x}{h}\\
  &=\sen x\cdot\lim_{h\to0}\frac{\cos h-1}{h}+\cos x\cdot\lim_{h\to0}\frac{\sen h}{h}\\
  &=\sen x\cdot0+\cos x\cdot1=\cos x .
\end{split}
\end{equation}
Allo stesso modo si calcola la derivata di $\cos x$, con la \eqref{eq:addizione-cos}:
\begin{equation}\label{eq:derivata-cos}
\begin{split}
  \lim_{h\to0}\frac{\cos(x+h)-\cos x}{h}
  &=\lim_{h\to0}\frac{\cos x\cos h-\sen x\sen h-\cos x}{h}\\
  &=\cos x\cdot\lim_{h\to0}\frac{\cos h-1}{h}-\sen x\cdot\lim_{h\to0}\frac{\sen h}{h}\\
  &=\cos x\cdot0-\sen x\cdot1=-\sen x .
\end{split}
\end{equation}
La derivata della funzione $\tg x$ si calcola con la regola di derivazione del quoziente,
nei punti in cui $\cos x\neq0$, usando la relazione fondamentale
\eqref{eq:relazione-fondamentale}:
\begin{equation}\label{eq:derivata-tg-calcolo}
\begin{split}
  D\tg x=D\Bigl(\frac{\sen x}{\cos x}\Bigr)&=\frac{D(\sen x)\cos x-\sen x\,D(\cos x)}{\cos^{2}x}\\
  &=\frac{\cos^{2}x+\sen^{2}x}{\cos^{2}x}=\frac1{\cos^{2}x};
\end{split}
\end{equation}
dividendo termine a termine si ottiene anche la forma $1+\tg^{2}x$.
\end{dimostrazione}

\begin{strategia}[Un'altra via]
Giusti \cite{giusti} trasforma la differenza $\sen(x+h)-\sen x$ in un prodotto con le
formule di prostaferesi: l'incremento $h$ compare allora in un solo fattore,
$\sen\frac h2$, e l'altro fattore è il coseno del punto medio. Basta così un solo limite
notevole, $\frac{\sen t}t\to1$, insieme alla continuità del coseno; il limite di
$\frac{1-\cos h}h$ non serve.
\end{strategia}

\begin{dimostrazione}[Seconda dimostrazione]
Per ogni $a$, $b$ reali, sottraendo membro a membro le formule di addizione
\eqref{eq:addizione-sen}, \eqref{eq:addizione-cos} scritte per $a+b$ e per $a-b$, si
ottengono le formule di prostaferesi (già usate nell'\cref{es:seno-lipschitz})
\begin{equation}\label{eq:prostaferesi-ab}
\begin{gathered}
  \sen(a+b)-\sen(a-b)=2\cos a\,\sen b,\\
  \cos(a+b)-\cos(a-b)=-2\sen a\,\sen b .
\end{gathered}
\end{equation}
Con $a=x+\frac h2$ e $b=\frac h2$ la prima dà, per $h\neq0$,
\[
  \frac{\sen(x+h)-\sen x}{h}=\cos\Bigl(x+\frac h2\Bigr)\cdot\frac{\sen(h/2)}{h/2}.
\]
Per $h\to0$ il primo fattore tende a $\cos x$, per la continuità del coseno
(\cref{prop:sen-cos-continue}); il secondo tende a $1$ per la \eqref{eq:lim-senx-x} e il
\cref{teo:limite-composta}, con la funzione interna $t=\frac h2$, diversa da zero per
$h\neq0$. Quindi $D\sen x=\cos x$. Allo stesso modo, per la seconda delle
\eqref{eq:prostaferesi-ab} e per la continuità del seno,
\[
  \frac{\cos(x+h)-\cos x}{h}=-\sen\Bigl(x+\frac h2\Bigr)\cdot\frac{\sen(h/2)}{h/2}
  \to-\sen x . \qedhere
\]
\end{dimostrazione}

\begin{osservazione}[La differenza centrata del seno]\label{oss:seno-centrata}
Con $a=x$ e $b=h$ la prima delle \eqref{eq:prostaferesi-ab} dà, per la differenza
centrata del \cref{sec:def-derivata},
\[
  \frac{\sen(x+h)-\sen(x-h)}{2h}=\cos x\cdot\frac{\sen h}{h}:
\]
la differenza centrata del seno è il valore esatto $\cos x$ moltiplicato per un fattore
che tende a $1$, e il suo errore vale esattamente $\cos x\,\bigl(1-\frac{\sen h}h\bigr)$
se $\cos x>0$. Per $x=1$ e $h=0{,}1$ si trova $0{,}5403\cdot0{,}0016658\approx9{,}001\cdot
10^{-4}$, il primo valore della colonna «errore centrato» della tabella del
\cref{sec:def-derivata}. Vedremo con la formula di Taylor che $1-\frac{\sen h}h$ è circa
$\frac{h^{2}}6$: per questo quell'errore si divide per $100$ quando $h$ si divide per $10$.
\end{osservazione}

\begin{attenzione}[Gli angoli vanno misurati in radianti]
Le formule \eqref{eq:derivata-sen-cos} valgono perché $\frac{\sen h}{h}\to1$, e questo
limite vale soltanto se l'angolo $h$ è misurato in radianti. Se l'angolo $\theta$ è
espresso in gradi, si ha $\sen\theta^{\circ}=\sen\frac{\pi\theta}{180}$ e, per la regola
della catena, $\frac{d}{d\theta}\sen\theta^{\circ}=\frac{\pi}{180}\cos\theta^{\circ}$: un
fattore $\frac\pi{180}\approx0{,}01745$ che nei calcoli di stabilità e di prestazioni è
fonte di errori classici (ad esempio quando si usa una pendenza della retta di portanza
$C_{L\alpha}$ espressa «per grado» in una formula scritta «per radiante»).
\end{attenzione}

Riassumiamo nella \cref{tab:derivate-elementari} le principali formule di derivazione
trovate in questo paragrafo.

\begin{table}[htb]
\centering
\renewcommand{\arraystretch}{1.45}
\begin{tabular}{l l l}
\toprule
$f(x)$ & $f'(x)$ & validità\\
\midrule
$x^{b}$ & $bx^{b-1}$ & $x>0$ ($x\in\R$ se $b\in\N$)\\
$\log x$ & $\dfrac1x$ & $x>0$\\
$\log_{a}x$ & $\dfrac1{x\log a}$ & $x>0$\\
$e^{x}$ & $e^{x}$ & $x\in\R$\\
$a^{x}$ & $a^{x}\log a$ & $x\in\R$\\
$\sen x$ & $\cos x$ & $x\in\R$\\
$\cos x$ & $-\sen x$ & $x\in\R$\\
$\tg x$ & $\dfrac1{\cos^{2}x}=1+\tg^{2}x$ & $x\neq\frac\pi2+k\pi$\\
\bottomrule
\end{tabular}
\caption{Le derivate delle funzioni elementari. Con le regole dei
\cref{sec:operazioni-derivate,sec:composte-inverse} si derivano tutte le funzioni che da
esse si ottengono con le operazioni algebriche, la composizione e l'inversione.}
\label{tab:derivate-elementari}
\end{table}

Con la tabella e le regole di calcolo il lettore può ormai derivare qualsiasi funzione
elementare. Un sistema di calcolo simbolico come \texttt{sympy} applica esattamente le
stesse regole, ed è utile per controllare i conti.

\begin{python}[Derivate calcolate con \texttt{sympy}]
import sympy as sp

x = sp.symbols("x", positive=True)
funzioni = [x**3 * sp.sin(x), sp.sin(x**2), sp.sin(x)**2, sp.tan(x),
            sp.exp(x) / x, sp.log(x**2 + 1), x**x, sp.atan(1 / x)]
for f in funzioni:
    print(f"D[{f}] = {sp.simplify(sp.diff(f, x))}")
\end{python}
\begin{risultato}
D[x**3*sin(x)] = x**2*(x*cos(x) + 3*sin(x))
D[sin(x**2)] = 2*x*cos(x**2)
D[sin(x)**2] = sin(2*x)
D[tan(x)] = cos(x)**(-2)
D[exp(x)/x] = (x - 1)*exp(x)/x**2
D[log(x**2 + 1)] = 2*x/(x**2 + 1)
D[x**x] = x**x*(log(x) + 1)
D[atan(1/x)] = -1/(x**2 + 1)
\end{risultato}

Il lettore verifichi a mano ciascun risultato. Si noti che \texttt{sympy} scrive
$\sen2x$ al posto di $2\sen x\cos x$ (formula di duplicazione) e che l'ultima riga contiene
la derivata di $\operatorname{arctg}\frac1x$, che si calcola con quella dell'arcotangente
(\cref{sec:trig-inverse}): è opposta alla derivata di $\operatorname{arctg}x$, il che non è
un caso (\cref{es:arctg-1-su-x}).

\begin{idea}[Come varia la densità dell'aria durante una salita]
Nella troposfera dell'atmosfera standard (ISA) la densità è data, in funzione della quota
(geopotenziale) $h$ in metri, da
\begin{gather*}
  \rho(h)=\rho_{0}\Bigl(1-\frac{L\,h}{T_{0}}\Bigr)^{n},\\
  \rho_{0}=1{,}225\ \text{kg/m}^{3},\quad T_{0}=288{,}15\ \text{K},\quad L=0{,}0065\ \text{K/m},
\end{gather*}
con $n=\frac{g}{RL}-1\approx4{,}2559$, dove $g=9{,}80665$~m/s$^{2}$ è l'accelerazione di
gravità standard e $R=287{,}05$~J/(kg\,K) la costante dei gas per l'aria. Per la regola
della catena e la \eqref{eq:derivata-xb},
\[
  \frac{d\rho}{dh}=\rho_{0}\,n\Bigl(1-\frac{Lh}{T_{0}}\Bigr)^{n-1}\cdot\Bigl(-\frac L{T_{0}}\Bigr),
\]
e, se il velivolo sale con velocità verticale $v_{z}=\frac{dh}{dt}$, ancora per la regola
della catena la densità varia nel tempo con tasso $\frac{d\rho}{dt}=\frac{d\rho}{dh}\,v_{z}$.
Il programma che segue calcola questi valori per $v_{z}=5$~m/s, controllando la formula
con una differenza centrata.
\end{idea}

\begin{python}[Derivata della densità ISA rispetto alla quota e al tempo]
# Densità ISA (troposfera): rho0 (1 - L h/T0)^n, n = g/(R L) - 1
rho0, T0, L, g, R = 1.225, 288.15, 0.0065, 9.80665, 287.05
n = g / (R * L) - 1
rho = lambda h: rho0 * (1 - L * h / T0) ** n
# derivata con la regola della catena: D[u^n] = n u^(n-1) u'
drho = lambda h: rho0 * n * (1 - L * h / T0) ** (n - 1) * (-L / T0)

vz = 5.0                               # velocità verticale (m/s)
print(f"n = {n:.4f}")
print("    h      rho      drho/dh     numerica     drho/dt")
for h in (0, 2000, 5000, 10000):
    num = (rho(h + 0.01) - rho(h - 0.01)) / 0.02  # diff. centrata
    print(f"{h:5d}   {rho(h):.4f}  {drho(h):.4e}  {num:.4e}  "
          f"{drho(h) * vz:.3e}")
\end{python}
\begin{risultato}
n = 4.2559
    h      rho      drho/dh     numerica     drho/dt
    0   1.2250  -1.1760e-04  -1.1760e-04  -5.880e-04
 2000   1.0065  -1.0119e-04  -1.0119e-04  -5.060e-04
 5000   0.7361  -7.9654e-05  -7.9654e-05  -3.983e-04
10000   0.4127  -5.1162e-05  -5.1162e-05  -2.558e-04
\end{risultato}

Al livello del mare la densità diminuisce di circa $1{,}2\cdot10^{-4}$~kg/m$^{3}$ per
ogni metro di quota, e salendo a $5$~m/s di circa $5{,}9\cdot10^{-4}$~kg/m$^{3}$ al
secondo; a $10\,000$~m le variazioni sono meno della metà: dalla formula si ricava
$\frac{d\rho}{dh}=-\frac{nL}{T}\,\rho$, e a quella quota la densità è circa un terzo di quella
al suolo, mentre la temperatura $T=T_{0}-Lh$ è scesa soltanto del $23\%$ circa.
